\begin{table}[htbp]
\centering
\small
\caption{Сдвиг признаков при сокращении и расширении текста: снятая интерпретация через объём}
\label{tab:s15}
\begin{tabular}{>{\RaggedRight\hspace{0pt}\arraybackslash}p{0.302\linewidth}>{\RaggedRight\hspace{0pt}\arraybackslash}p{0.302\linewidth}>{\RaggedRight\hspace{0pt}\arraybackslash}p{0.302\linewidth}}
\toprule
Признак & t10 сокращение: медиана Δz (документов) & t11 расширение: медиана Δz (документов) \\
\midrule
M01 разброс сходства соседних предложений & +0.13 (51) & +7.86 (60) \\
S08 & −0.27 (46) & +4.53 (60) \\
L02 & +0.11 (42) & −3.02 (60) \\
L01 & +0.15 (44) & −1.45 (60) \\
R06 & −1.30 (23) & −1.39 (23) \\
R07 & −1.00 (22) & −1.00 (22) \\
P01 & −0.37 (24) & −0.34 (27) \\
F01 & −0.17 (15) & +0.10 (6) \\
\bottomrule
\end{tabular}
\begin{minipage}{\linewidth}\footnotesize\vspace{2pt}
\textit{Примечание.} Единица анализа — документ панели. Знаменатель: 60 документов панели; в скобках — число документов, у которых признак сдвинулся. Данные: стресс-ревизия r4, процедура 1; стандартизация по 1882 документам матрицы признаков v5. Медиана стандартизованного сдвига Δz по документам панели. Шкала: знак показывает направление сдвига признака, не решения. Сокращения: Δz — сдвиг в стандартных отклонениях признака; M01 — стандартное отклонение косинусов соседних пар предложений, а не среднее сходство. Пропуски: интерпретация «оба преобразования меняют объём, поэтому решение рушится» снята: форматная четвёрка сдвигается у обоих одинаково, а расходятся преобразования на M01. Поправка на множественные сравнения не применялась.
\end{minipage}
\end{table}
